\section{Uniform perturbation from exactness}\label{sec:exactness}

We now begin the exact proof that \propF\ implies \propP.  Its first
step is a perturbation estimate that applies simultaneously to every
primitive quotient in which a fixed positive element has norm at least
a prescribed positive threshold.  The constant in that estimate comes
from the following obstruction: there is one positive error threshold,
independent of the unital exact algebra and the matrix size,
below which a contraction cannot contain approximately reducing copies
of every finite matrix contraction.

\subsection{The universal contraction obstruction}

Recall that a $C^*$-algebra $B$ is \emph{exact} if, for every closed
two-sided ideal $I$ of every $C^*$-algebra $C$, the canonical map
\[
 (B\otimes_{\min}C)/(B\otimes_{\min}I)
 \longrightarrow B\otimes_{\min}(C/I)
\]
is an isomorphism.  Here and below, homomorphisms of $C^*$-algebras are
$*$-homomorphisms.  In the obstruction, the same finite-dimensional
copy must approximately intertwine both a contraction and its adjoint.

\begin{proposition}[Uniform gap]\label{prop:uniform-gap}
There is a constant $\gamma>0$ such that no unital exact $C^*$-algebra
$B$ contains a contraction $v$ with the following property: for every
integer $d\geq1$ and every contraction $D\in M_d(\C)$, there are a unital
representation $\pi:B\to B(H_\pi)$ and an isometry
$J:\C^d\to H_\pi$ satisfying
\begin{equation}\label{eq:forbidden-models}
 \|\pi(v)J-JD\|<\gamma,
 \qquad
 \|\pi(v)^*J-JD^*\|<\gamma.
\end{equation}
The representation and the isometry are allowed to depend on $d$ and
$D$.
\end{proposition}

The proof has two parts.  The universal unital algebra of a contraction
fails exactness, as witnessed by a tensor quotient.  If exact targets
admitted the proposed reducing copies with arbitrarily small common
error, their evaluations would give a lower bound on every matrix norm
of each fixed polynomial space.  Holding that space and one evaluation
fixed while an ideal approximate unit varies
would force the tensor quotient map to be injective.  The two relations
in \eqref{eq:forbidden-models} are what allow this argument to test words
containing both the generator and its adjoint.

Let $\U=C^*(1,u)$ be the universal
unital $C^*$-algebra generated by a contraction.  Thus every contraction
$T$ on a Hilbert space determines a unital homomorphism $u\mapsto T$.
One construction is to complete the unital $*$-algebra of polynomials
in $u,u^*$ in the supremum seminorm over all such evaluations, after
dividing out its null space.  The seminorm of each polynomial is finite
because a word in a contraction and its adjoint has norm at most one.

We shall use the following elementary quotient-norm formula.  If $D$
is unital and $(e_\lambda)$ is a positive contractive approximate unit
of an ideal $I\subset C$, then, for $z\in D\otimes_{\min}C$,
\begin{equation}\label{eq:tensor-quotient-norm}
 \|z+D\otimes_{\min}I\|
 =\lim_\lambda\|z(1\otimes(1-e_\lambda))\|.
\end{equation}
When $C$ is nonunital, the multiplication is performed in its
unitization; its result still belongs to $D\otimes_{\min}C$.
Indeed, $1\otimes e_\lambda$ is an approximate unit for
$D\otimes_{\min}I$.  The product on the right represents the same
quotient class as $z$, so its norm bounds the left side from below.
For any $y\in D\otimes_{\min}I$, contractivity gives
\[
 \limsup_\lambda\|z(1\otimes(1-e_\lambda))\|
 \leq \|z-y\|
       +\limsup_\lambda\|y(1\otimes(1-e_\lambda))\|
 =\|z-y\|.
\]
Taking the infimum over $y$ proves the formula, including existence of
the limit.

The formula also gives the form of subalgebra permanence that we need:
exactness passes to unital $C^*$-subalgebras; see also
\cite[IV.3.4.3]{Blackadar}.  For a unital inclusion $D\subset B$, the spatial
tensor norms on $D\otimes C$ and $D\otimes(C/I)$ are the norms
inherited from the corresponding tensor products with $B$.
Formula~\eqref{eq:tensor-quotient-norm} therefore gives, for
$z\in D\otimes_{\min}C$,
\[
 \|z+D\otimes_{\min}I\|
 =\|z+B\otimes_{\min}I\|.
\]
If $B$ is exact, this equals
$\|(\id_D\otimes q)z\|$, where $q:C\to C/I$ is the quotient map.
This proves exactness of $D$.

\begin{lemma}\label{lem:nonexact-universal}
The universal contraction algebra $\U$ is not exact.
\end{lemma}

\begin{proof}
We first embed the full group algebra $Q=C^*(\F_2)$ into $\U$.
The two free generators are sent to
\[
 \exp(2\pi i\operatorname{Re}u),
 \qquad \exp(2\pi i\operatorname{Im}u).
\]
To see that the resulting homomorphism is injective, take any pair of
unitaries $S,T$ on a Hilbert space.  Bounded Borel functional calculus
gives selfadjoint contractions $H,K$ such that
$S=\exp(\pi iH)$ and $T=\exp(\pi iK)$.  The contraction
$(H+iK)/2$ evaluates the two displayed unitaries to $S,T$.
Thus every unital representation of $Q$ is a restriction of an evaluation of
$\U$, and the homomorphism is isometric.  The logarithms need only
belong to the bounded operators on the representation space.
Courtney and Sherman give a related embedding with a different
normalization in \cite[Theorem~6.10]{CS}.

Nonexactness of $C^*(\F_2)$ was established by
Wassermann \cite{Wassermann1976}.  Here is an explicit proof for $Q$.
We construct a representation of $Q\otimes_{\min}Q$ that annihilates
$Q\otimes_{\min}I$, where $I$ is the kernel of the regular quotient.
Exactness would make it descend to $Q\otimes_{\min}(Q/I)$.  A concrete
tensor sum, however, has spatial norm in that algebra smaller than the
norm the proposed descended representation would assign to it.
This will exhibit the required failure of the tensor quotient map.
Put $G=\F_2$, and
write $u_s$ for its canonical unitaries in $Q$.  On $\ell^2(G)$ let
\[
 \lambda_s\delta_g=\delta_{sg},
 \qquad \rho_s\delta_g=\delta_{gs^{-1}}.
\]
There is a state $\varphi$ on $Q\otimes_{\min}Q$ with
\begin{equation}\label{eq:diagonal-state}
 \varphi(u_s\otimes u_t)=
 \begin{cases}1,&s=t,\\0,&s\ne t.\end{cases}
\end{equation}
For completeness, finite quotients supplying this state can be
constructed directly.  On the finite word ball of radius $m$ in $G$,
partially defined left translation by each of the two generators is
an injection between subsets of equal size, hence extends to a
permutation of the ball.  Freeness gives a homomorphism $q_m$ from
$G$ to the finite permutation group generated by these extensions.
The inverse of each extension agrees with inverse left translation
whenever both endpoints lie in the ball.
If a reduced word has length at most $m$, applying its permutations
to the identity follows its successive suffixes inside the ball and
ends at that word.  Consequently $q_m$ distinguishes distinct elements
of the word ball of radius $m$.

For the finite group $F_m=q_m(G)$, its commuting left and right
regular representations give a representation of
$C^*(F_m)\otimes_{\min}C^*(F_m)$: these algebras are finite direct
sums of matrix algebras, for which the spatial and maximal tensor
norms agree.  Pulling its vector state at the identity back to
$Q\otimes_{\min}Q$ gives the coefficient
$\boldsymbol{1}_{q_m(s)=q_m(t)}$ on $u_s\otimes u_t$.
These coefficients tend to those in \eqref{eq:diagonal-state}.
The group tensors span a dense $*$-subalgebra, and all the states have
norm one, so their limit on that span extends to a state on the
spatial completion.

The GNS representation of $\varphi$ is precisely
\begin{equation}\label{eq:left-right-gns}
 u_s\otimes u_t\longmapsto\lambda_s\rho_t
 \quad\hbox{on }\ell^2(G),
\end{equation}
with cyclic vector $\delta_e$.  Indeed the cyclic vector associated
with $u_s\otimes u_t$ corresponds to $\delta_{st^{-1}}$; the inner
products agree because
\[
 s^{-1}a=t^{-1}b\quad\Longleftrightarrow\quad st^{-1}=ab^{-1},
\]
and left multiplication by $u_x\otimes u_y$ becomes
$\lambda_x\rho_y$.  These vectors span $\ell^2(G)$.

Let $I=\ker(Q\to C_r^*(G))$.  The linear unitary
$R\delta_g=\delta_{g^{-1}}$ satisfies $R\lambda_sR^*=\rho_s$.
Thus the representation in \eqref{eq:left-right-gns} annihilates
$Q\otimes_{\min}I$.  Exactness of $Q$ would make it descend to a
representation of $Q\otimes_{\min}C_r^*(G)$, in which every
$u_s\otimes\lambda_s$ fixes $\delta_e$.

This contradicts a spatial norm bound.  In a faithful unital
representation $\pi:Q\to B(H)$, the unitary
\[
 V(\xi\otimes\delta_g)=\pi(u_g)\xi\otimes\delta_g
\]
satisfies
$V(1\otimes\lambda_s)V^*=\pi(u_s)\otimes\lambda_s$.
If $a,b$ are the free generators, the norm of
\[
 u_a\otimes\lambda_a+u_{a^{-1}}\otimes\lambda_{a^{-1}}
 +u_b\otimes\lambda_b+u_{b^{-1}}\otimes\lambda_{b^{-1}}
\]
in this faithful spatial representation is therefore the norm of the
adjacency operator on the $4$-regular tree.  That norm is at most
$2\sqrt3<4$.  To check the bound, orient each edge away from the
identity, and for a finitely supported vector $\xi$ bound the
contribution of a parent--child edge by
\[
 2|\xi(p)\xi(c)|
 \leq 3^{-1/2}|\xi(p)|^2+3^{1/2}|\xi(c)|^2.
\]
Every nonroot vertex has three children and one parent, and hence
total coefficient $2\sqrt3$; the root has coefficient
$4/\sqrt3\leq2\sqrt3$.  The quadratic-form bound extends by density
and proves the asserted operator norm bound.  In the purported
descended representation the same sum has eigenvalue $4$ on
$\delta_e$, which is impossible for a $*$-representation.
Thus $Q$ is not exact.  Its unital inclusion in $\U$ and the
inheritance proved above show that $\U$ is not exact either.
\end{proof}

We have obtained the algebraic obstruction.  To turn it into a fixed
error bound, we next test polynomial norms on finite contractions and
control those tests at every matrix amplification.  The latter
uniformity will allow one polynomial model to be held fixed while an
ideal approximate unit varies.

\begin{lemma}[Finite matrix tests]\label{lem:finite-matrix-tests}
For every integer $r\geq1$ and every $Z\in M_r(\U)$ whose entries are
$*$-polynomials in $u,u^*$,
\begin{equation}\label{eq:finite-matrix-tests}
 \|Z\|_{M_r(\U)}
 =\sup_{d\geq1,\ D\in M_d(\C),\ \|D\|\leq1}\|Z(D,D^*)\|.
\end{equation}
\end{lemma}

\begin{proof}
Contractivity of evaluations proves one inequality.  Every unital
representation of $\U$ is evaluation at the image of $u$, so these
evaluations also compute the matrix norm on the left.  For the other
inequality, evaluate at a contraction $T\in B(H)$, and take a unit vector
$\xi=(\xi_1,\ldots,\xi_r)\in H^r$ arbitrarily close to attaining
the norm of $Z(T,T^*)$.  If the entries of $Z$ have degree at most
$m$, let
\[
 H_0=\operatorname{span}
 \{w(T,T^*)\xi_j:|w|\leq m,\ 1\leq j\leq r\},
\]
where the empty word is included.  This is a finite-dimensional
nonzero space.  Set $D=P_{H_0}T|_{H_0}$, so that
$D^*=P_{H_0}T^*|_{H_0}$.  Applying a word successively from the right
never changes the relevant vectors under compression, since all its
suffix vectors lie in $H_0$.  Hence
\[
 w(D,D^*)\xi_j=w(T,T^*)\xi_j\quad (|w|\leq m),
\]
and $Z(D,D^*)\xi=Z(T,T^*)\xi$.  Taking the supremum over evaluations
and over these vectors proves \eqref{eq:finite-matrix-tests}.
The scalar norm-test formula is proved in \cite[Proposition~2.2]{CS};
the finite-word compression used here is the argument in the proof of
\cite[Theorem~3.1]{CS}. Including all components of the test vector
gives the matrix-valued assertion needed here.
\end{proof}

\begin{lemma}[A bound at every amplification]\label{lem:uniform-matrix-bound}
Let $E\subset\U$ be a finite-dimensional linear space consisting of
$*$-polynomials in $u,u^*$.  There is a finite constant $C_E\geq0$
with the following property.  Suppose $v$ is a contraction in a unital
$C^*$-algebra $B$, and fix one number $\theta>0$.  Assume that for every
integer $d\geq1$ and every contraction $D\in M_d(\C)$, there are a unital
representation $\pi:B\to B(H_\pi)$ and an isometry $J:\C^d\to H_\pi$
satisfying the two inequalities in \eqref{eq:forbidden-models} with
$\theta$ in place of $\gamma$.
For the homomorphism $\Phi:\U\to B$ with $\Phi(u)=v$, one then has
\begin{equation}\label{eq:all-amplification-bound}
 \|\Phi^{(r)}(Z)\|\geq(1-C_E\theta)\|Z\|
 \qquad (r\geq1,\ Z\in M_r(E)).
\end{equation}
\end{lemma}

\begin{proof}
If $E=0$, take $C_E=0$.  Otherwise choose a basis $P_1,\ldots,P_k$
of $E$, with coordinate functionals $f_1,\ldots,f_k$, and choose
polynomial representatives
$P_i=\sum_w a_{iw}w$.  Define
\begin{equation}\label{eq:explicit-CE}
 C_E=\sum_{i=1}^k\|f_i\|\sum_w|a_{iw}|\,|w|.
\end{equation}
Constants are allowed in the representatives; the empty word has
length zero.  For $Z=\sum_i B_i\otimes P_i\in M_r(E)$ and scalar
unit vectors $\alpha,\beta\in\C^r$,
\[
 |\alpha^*B_i\beta|
 =|f_i(\alpha^*Z\beta)|\leq\|f_i\|\|Z\|.
\]
It follows that $\|B_i\|\leq\|f_i\|\|Z\|$ for every matrix size $r$.

Given a finite contraction $D$, choose $\pi,J$ as in the hypothesis.
Both $\pi(v)$ and $D$ are contractions, and telescoping the two
intertwining relations gives
\[
 \|w(\pi(v),\pi(v)^*)J-Jw(D,D^*)\|\leq|w|\theta
\]
for every word $w$, including words in both letters.  More explicitly,
the difference for a word $A_1\cdots A_m$ is the sum of $m$ terms,
each consisting of contractions on either side of one difference
$\pi(v)J-JD$ or $\pi(v)^*J-JD^*$.  Put $J_r=1_{\C^r}\otimes J$.
The coefficient bound and \eqref{eq:explicit-CE} yield
\[
 \|\pi^{(r)}(\Phi^{(r)}(Z))J_r-J_r Z(D,D^*)\|
 \leq C_E\theta\|Z\|.
\]
Since $J_r$ is an isometry and $\pi^{(r)}$ is contractive,
\[
 \|Z(D,D^*)\|\leq\|\Phi^{(r)}(Z)\|+C_E\theta\|Z\|.
\]
Taking the supremum supplied by Lemma~\ref{lem:finite-matrix-tests}
proves \eqref{eq:all-amplification-bound}.  In particular, $C_E$
depends on the fixed polynomial space and its chosen representatives,
but on none of $r,d,D,B,v,\pi,J$.
\end{proof}

\begin{proof}[Proof of Proposition~\ref{prop:uniform-gap}]
If the Proposition were false, for every integer $j\geq1$ there
would be a unital exact algebra $B_j$ and a contraction $v_j\in B_j$
satisfying \eqref{eq:forbidden-models} with $1/j$ in place of $\gamma$.
Let $\Phi_j:\U\to B_j$ send $u$ to $v_j$.
Fix a finite-dimensional polynomial space $E$.  Before choosing any
matrix amplification, tensor factor, compression, or ideal, choose one
integer $j=j(E)\geq\max\{1,2C_E\}$.  Lemma~\ref{lem:uniform-matrix-bound}
then gives
\begin{equation}\label{eq:fixed-j-matrix-bound}
 \|Z\|\leq2\|\Phi_j^{(r)}(Z)\|
 \qquad(r\geq1,\ Z\in M_r(E)).
\end{equation}

This same $j$ works after tensoring with an arbitrary $C^*$-algebra
$C$: for every $z\in E\odot C$,
\begin{equation}\label{eq:fixed-j-tensor-bound}
 \|z\|\leq2\|(\Phi_j\otimes\id_C)z\|,
\end{equation}
where $\odot$ denotes the algebraic tensor product and all norms are
spatial.  The assertion is immediate if $C=0$.  Otherwise, represent
$C$ faithfully on a nonzero Hilbert space $K$, and use
faithful representations of $\U$ and $B_j$ for their respective
spatial tensor products.  Compress the $C$ representation to a
nonzero finite-dimensional subspace $K_0\subset K$.  If
$z=\sum_i P_i\otimes c_i$, its compression has coefficients
$P_{K_0}c_i|_{K_0}\in M_d(\C)$ and, after interchanging the tensor
factors, belongs to $M_d(E)$.  Inequality~\eqref{eq:fixed-j-matrix-bound}
bounds its norm by twice the norm of the corresponding compression
on the $B_j$ side, which is at most
$2\|(\Phi_j\otimes\id_C)z\|$.  The norms of these compressions
have supremum $\|z\|$: as $K_0$ increases, the projections
$1\otimes P_{K_0}$ converge strongly to the identity, and their
two-sided compressions converge strongly to $z$.  This proves
\eqref{eq:fixed-j-tensor-bound}.

Now fix any ideal $I\subset C$, with quotient map $q:C\to C/I$
and a positive contractive approximate unit $(e_\lambda)$ of $I$.
For $z\in E\odot C$, every
$z(1\otimes(1-e_\lambda))$ still belongs to $E\odot C$.
Consequently \eqref{eq:fixed-j-tensor-bound} applies to the entire
net with the single index $j(E)$.  Formula~\eqref{eq:tensor-quotient-norm}
and exactness of $B_j$ give
\begin{align*}
 \|z+\U\otimes_{\min}I\|
 &=\lim_\lambda\|z(1\otimes(1-e_\lambda))\|\\
 &\leq2\lim_\lambda
   \|(\Phi_j\otimes\id_C)(z)(1\otimes(1-e_\lambda))\|\\
 &=2\|(\Phi_j\otimes\id_{C/I})((\id_\U\otimes q)z)\|\\
 &\leq2\|(\id_\U\otimes q)z\|.
\end{align*}
Finite sums of polynomial tensors are dense in $\U\otimes_{\min}C$.
Although $j(E)$ may vary with the polynomial space, the resulting
constant $2$ does not.  The last inequality therefore extends by
continuity to every $z\in\U\otimes_{\min}C$.  It says that the
canonical surjective homomorphism
\[
 \kappa:(\U\otimes_{\min}C)/(\U\otimes_{\min}I)
 \longrightarrow\U\otimes_{\min}(C/I)
\]
has zero kernel: if $\kappa(z+\U\otimes_{\min}I)=0$, the quotient
class has norm zero.  An injective $*$-homomorphism of $C^*$-algebras
is isometric, so the fixed factor $2$ suffices to prove that $\kappa$
is an isomorphism.  Since $C,I$ were arbitrary, $\U$ would be exact,
contrary to Lemma~\ref{lem:nonexact-universal}.
\end{proof}

Proposition~\ref{prop:uniform-gap} therefore supplies a single
$\gamma>0$ for all unital exact targets and all matrix sizes.  The remaining
part of this section converts that obstruction into a perturbation
with a quotient norm lower bound.
